\section{서론}
\label{sec:intro}

\begin{revision}

비전--언어 주행 모델은 장면을 해석하고 행동 이유를 언어로 설명한다 \cite{sima2024drivelm,tian2024drivevlm}. Chain-of-Causation(CoC)은 이런 행동 이유를 인과관계로 연결한 주석이다. 하지만 이 방식은 같은 장면의 CoC들이 개별적으로는 자연스러워도 시간순으로 함께 읽으면 양립하지 않을 수 있다는 문제점이 있다. 생성ㆍ작성 과정에서 조건이 누락되거나, 앞서 제시한 요구와 양립하지 않는 행동이 기록될 수 있기 때문이다. 이러한 불일치 때문에 학습 자료의 행동 요구가 모호해지고, 같은 자료를 다시 검사했을 때 일관된 결과를 얻기 어렵다.

다음은 연구 동기를 설명하는 가상 예시이며 자연 자료에서 발견한 오류가 아니다. 앞선 CoC가 ``보행자가 통과할 때까지 정지한다''고 하고, 후속 CoC가 ``신호가 녹색으로 바뀌었으므로 출발한다''고 했다고 하자. 녹색 신호는 출발의 새로운 근거이지만 보행자가 통과했다는 증거는 아니다. 보행자가 여전히 횡단 중임이 확인되면 출발은 이전 정지 요구와 충돌한다. 통과 여부가 제시되지 않았다면 확정 모순이 아니라 정지 요구의 해제 근거가 부족한 경우이다. 보행자의 통과가 확인된 뒤 출발하면 해당 요구와 양립한다. 새로운 행동 근거 B가 있다는 이유만으로 이전 요구 A를 지워서는 안 된다.

본 연구는 이러한 시간 순서와 조건의 관계를 UPPAAL 모델로 검사하고 문제 사건과 근거를 추적한다. UPPAAL은 Uppsala 대학과 Aalborg 대학이 공동 개발한 시간 오토마타 기반 모델 검사 도구의 이름이며, 두 대학 이름을 연결한 명칭이다 \cite{uppaalofficial,uppaalname,behrmann2004uppaal}. 문장의 요구를 구조화하는 과정과 그 조건의 참ㆍ거짓을 확인하는 과정은 별개다. 현재 자동 변환기는 제한된 영어 행동ㆍ조건 문구를 찾지만 관측 증거를 확정하지 않는다. 따라서 자동 텍스트 경로와 설계자가 작성한 조건별 프레임 경로를 구분한다.

\begin{figure*}[t]
\centering
\color{revisionblue}
\begin{tikzpicture}[
 box/.style={draw=revisionblue,rounded corners,align=center,text width=3.4cm,minimum height=11mm,font=\footnotesize,text=revisionblue},
 flow/.style={-{Latex},draw=revisionblue,thick}]
\node[box] (text) at (0,0) {Stored CoC text\\rule-based phrase extraction\\observed evidence unresolved};
\node[box] (frame) at (0,-1.7) {Authored text + frames\\condition IDs and source spans\\fixed scenario expectations};
\node[box] (check) at (5,-.85) {Frozen single-obligation path\\FSM comparator\\condition-indexed UPPAAL};
\node[box] (trace) at (10,-.85) {Decision + first issue event\\obligation + reason\\XML / queries / logs};
\draw[flow] (text.east) -- (check.north west);
\draw[flow] (frame.east) -- (check.south west);
\draw[flow] (check.east) -- (trace.west);
\end{tikzpicture}

\bifigcaption{자동 문구 추출과 작성 프레임을 구분한 검사 흐름. 원문 구절, 조건별 의무, 최초 문제 사건과 사유를 연결한다.}{Separate automatic-text and authored-frame paths with links from source spans to obligations and issue records.}
\label{fig:overview-pipeline}
\end{figure*}

연구 질문은 ``연속 CoC의 의무와 해제 근거를 연결하여 행동 전환의 모순과 근거 부족을 구분하고, 최초 문제 사건ㆍ관련 의무ㆍ사유를 추적할 수 있는가?''이다. 이를 확인하기 위해 기존 단일 의무 검사기의 최초 결과를 보존하고, 이전 사건 1/3/5개를 보는 검사와 별도 유한상태기계(FSM)를 비교했다. 최초 실패 분석 후 조건별 UPPAAL 모델을 개발하여 같은 사례를 재시험했다. 이 재시험을 독립 성능 검증으로 해석하지 않는다.

기여는 세 가지다. 첫째, 원문 구절--의무 ID--해제 조건--후속 증거--행동을 연결하고 모순과 근거 부족을 구분하는 표현을 제시한다. 둘째, 조건별 상태와 속성을 UPPAAL로 실행하고 최초 문제 위치와 사유를 추적하는 절차를 구체화한다. 셋째, 최초 통제 시험, 실패 분석, 동일 사례 재시험과 자연 자료의 해석 한계를 구분하여 적용 가능성을 보고한다. 단순 FSM도 같은 사례를 모두 판정했으므로 UPPAAL의 정확도 우월성을 주장하지 않는다.

이하에서는 자료와 조건 표현을 정의한 뒤, 원문에서 프레임ㆍ모델ㆍ검토 기록으로 이어지는 절차를 설명한다. 최초 시험과 실패 분석 후 재시험을 구분하여 결과를 제시하고, 주석 검토에서의 활용과 남은 검증 범위를 논의한다.
\end{revision}
